\documentclass[11pt]{article}
\usepackage{times}
\usepackage{fullpage}
\usepackage{url}
\usepackage{hyperref}
\usepackage{fancyhdr}
\usepackage{graphicx}
\usepackage{enumitem}
\usepackage{subcaption}
\usepackage{amsmath, amsfonts, amsthm, fouriernc}
\usepackage{IEEEtrantools}
\usepackage{parskip}

\pagestyle{fancy}
\addtolength{\headheight}{2em}
\addtolength{\headsep}{1.5em}
\lhead{\large{ME 2801}}

\usepackage{sectsty}
\sectionfont{\fontsize{12}{8}\selectfont}

\begin{document}

\begin{center}
  \Large{\bf{Assignment: Time Response}}
\end{center}

The textbook exercises, designated with the word ``Nise'', are from the
``Problems'' section at the end of Chapter~4 in \emph{Control Systems
Engineering} by Norman Nise, 7th edition.

This assignment uses exercises from the textbook.  The additional instructions are to clarify and/or expand each exercise.

% ---------------------------------------------------------------
\section*{Nise 4.2 and 4.3: Step response of first-order system}

For the MATLAB plots, add annotation (by hand or using the computer) to indicate
the \emph{time constant}~($\tau$) of the system.

% ---------------------------------------------------------------
\section*{Nise 4.8: Pole-zero maps and step responses}

\begin{itemize}
  \item Generate the pole-zero plots by hand; a simple sketch of the $s$-plane
    is sufficient.
    \begin{itemize}
      \item The MATLAB function \texttt{roots()} finds roots of polynomials
        (useful if you want to avoid the quadratic formula).
      \item Optional: \texttt{pzmap()} checks your work directly from the
        transfer function.
    \end{itemize}
  \item Use the MATLAB \texttt{step()} command to plot the step response.
    The goal is to build intuition connecting pole-zero locations to
    time-response character.
  \item The \emph{general forms} referred to in the problem statement are:
    \begin{IEEEeqnarray*}{rCl}
      c(t) & = & K[1-e^{-at}]  \\
      c(t) & = & K[1-Ae^{-\sigma_1 t}-Be^{-\sigma_2 t}] \\
      c(t) & = & K[1-Ae^{-\sigma_d t} \cos(\omega_d t - \phi)] \\
      c(t) & = & K[1 - A \cos(\omega_1 t - \phi)] \\
      c(t) & = & K[1 - Ae^{-\sigma_1 t} - B\,t\,e^{-\sigma_1 t}]
    \end{IEEEeqnarray*}
    Without computing the full inverse Laplace transform, identify numeric
    values for $a$, $\sigma_1$, $\sigma_2$, $\sigma_d$, $\omega_d$, $\omega_1$
    from the appropriate general form.  You do not need to find the coefficients
    $K$, $A$, $B$ (they are both tedious to compute by hand and not necessary for understanding the time-response characteristics).
\end{itemize}

% ---------------------------------------------------------------
\section*{Nise 4.20 and 4.21: Response of second-order system}

Parts~(a) and~(b) only.

\subsection*{Nise 4.20: Computing performance metrics}

\subsubsection*{Rise Time Approximation: Second-Order Underdamped System}
For a second-order underdamped system with transfer function
\[
  G(s) = \frac{\omega_n^2}{s^2 + 2\zeta\omega_n s + \omega_n^2},
\]
``a precise analytical relationship between rise time and damping ratio, $\zeta$, cannot be found''~\cite[Sec.~4.6]{nise2015}.  Consequently, we need a method to approximate the rise time.  

The Nise textbook offers both a polynomial approximation and a graphical method for estimating the rise time.  The polynomial approximation from
the textbook (footnote~5 to Section~4.6):
\[
  \omega_n T_r \approx 1.76\zeta^3 - 0.417\zeta^2 + 1.039\zeta + 1
  \qquad (0 < \zeta < 0.9,\;\text{error} < 0.5\%)
\]
Alternatively, Figure~4.16 in the textbook can be used as a graphical lookup table for estimating the rise time.


An alternative coarse (but useful) approximation from Franklin, Powell, and Emami-Naeini~\cite[Sec.~3.4]{franklin2010}, derived for $\zeta \approx 0.5$, is
\[
T_r \approx \frac{1.8}{\omega_n}
\]
For design work, where the relative rise-time performance between different designs is more important than the exact value, this coarse approximation can be quite useful.

For the purposes of class, the coarse approximation is sufficient, but students are free to use the more precise polynomial approximation if desired.  Do remember that these approximations only hold for underdamped second-order systems in the form shown above.  


\section*{Nise 4.32: Step response from pole-zero description}

Check your answers by using MATLAB to plot the step response for your answers.  Verify that
the generated plots agree with the graphs in the exercise.

\begin{thebibliography}{9}
\bibitem{nise2015} N.~S. Nise, \emph{Control Systems Engineering}, 7th ed. Hoboken, NJ: Wiley, 2015.
\bibitem{franklin2010} G.~F. Franklin, J.~D. Powell, and A. Emami-Naeini, \emph{Feedback Control of Dynamic Systems}, 6th ed. Upper Saddle River, NJ: Pearson, 2010.
\end{thebibliography}

\end{document}
